\section{控制算法设计}

\subsection{动力学补偿与误差通道规范化}

对等效关节空间模型 \eqref{eq:joint_model}，采用如下动力学补偿结构
\begin{equation}
\tau=M_e(q_m)v+C_e(q_m,\dot q_m)\dot q_m-\hat d(t),
\label{eq:computed_torque}
\end{equation}
其中 \(v\) 为辅助加速度输入，\(\hat d(t)\) 为扰动补偿项。将式 \eqref{eq:computed_torque} 代入式 \eqref{eq:joint_model} 可得
\[
M_e(q_m)(\ddot q_m-v)=d(t)-\hat d(t).
\]
在理想补偿情形下 \(\hat d(t)=d(t)\)，且 \(M_e(q_m)\) 正定可逆，因此
\[
\ddot q_m=v.
\]

令
\[
v=\ddot q_d+\nu ,
\]
并定义
\[
e=q_m-q_d,\qquad
z=\begin{bmatrix}e\\ \dot e\end{bmatrix}.
\]
则
\[
\ddot e=\ddot q_m-\ddot q_d=\nu,
\]
误差系统可写成二阶积分链形式
\begin{equation}
\dot z=Az+B\nu ,
\label{eq:linear_error}
\end{equation}
其中
\[
A=\begin{bmatrix}0&I_n\\0&0\end{bmatrix},\qquad
B=\begin{bmatrix}0\\I_n\end{bmatrix}.
\]
矩阵对 \((A,B)\) 可控，因此 PDF 增益的设计对象由原非线性动力学转移为线性可控误差通道。这个步骤是本文方法成立的关键：PDF 理论只需作用在 \eqref{eq:linear_error} 上，而不直接处理自由漂浮机械臂的非线性耦合项。

\subsection{PDF 线性误差系统设计}

文献\cite{zhouFixedtimeStabilizationLinear2022}考虑的基本对象是可控线性时滞系统，其固定时间镇定依赖如下周期延迟结构：
\begin{equation}
\dot x(t)=F(t)x(t)+G(t)x(t-h),
\label{eq:periodic_delay_general}
\end{equation}
其中 \(F(t),G(t)\) 为 \(2h\) 周期矩阵，且 \(G(t)=0\) 在每个周期的前半段 \([0,h]\) 成立。由分步法可知，系统在一个周期后的状态满足
\[
x(2(k+1)h)=\Delta(h)x(2kh),
\]
其中 monodromy matrix 为
\begin{equation}
\Delta(h)=\Phi(2h,0)+\int_h^{2h}\Phi(2h,s)G_0(s)\Phi(s-h,0)\,ds.
\label{eq:monodromy}
\end{equation}
若 \(\Delta(h)\) 为幂零矩阵，即存在最小整数 \(\nu_0\ge1\) 使 \(\Delta^{\nu_0}(h)=0\)，则 \(x(t)=0,\ t\ge 2\nu_0h\)。特别地，若通过增益设计使 \(\Delta(h)=0\)，则系统可在一个完整周期后归零\cite{zhouFixedtimeStabilizationLinear2022}。

为将该思想用于式 \eqref{eq:linear_error}，先取常值反馈 \(K_0\) 使
\[
A_c=A+BK_0
\]
成为期望的基准闭环矩阵，再引入周期延迟反馈
\begin{equation}
\nu(t)=K_0z(t)-K_c(t)z(t-h).
\label{eq:pdf_error_input}
\end{equation}
此时误差闭环为
\begin{equation}
\dot z(t)=A_cz(t)-BK_c(t)z(t-h).
\label{eq:closed_delay_error}
\end{equation}
令 \(R_h(t)\) 为 \(2h\) 周期光滑函数，满足
\[
R_h(t)=0,\quad t\in[0,h],\qquad
R_h(t)\ge0,\quad t\in[h,2h],
\]
且在 \([h,2h]\) 的某个子区间上正定。根据文献\cite{zhouFixedtimeStabilizationLinear2022}的构造，可定义
\begin{equation}
W_c(A_c,h)=\int_h^{2h}e^{-A_cs}BR_h(s)B^\top e^{-A_c^\top s}\,ds,
\label{eq:wc_matrix}
\end{equation}
并取
\begin{equation}
K_c(t)=R_h(t)B^\top e^{-A_c^\top t}W_c^{-1}(A_c,h)e^{A_c(h-t)}.
\label{eq:kc_pdf}
\end{equation}
由于 \((A_c,B)\) 与 \((A,B)\) 同可控，且 \(R_h(t)\) 在后半周期内正定，\(W_c(A_c,h)\) 非奇异。将 \eqref{eq:kc_pdf} 代入 \eqref{eq:monodromy} 可使闭环的 monodromy matrix 满足 \(\Delta(h)=0\)，从而得到固定时间零化性质。式 \eqref{eq:kc_pdf} 也说明 PDF 增益可通过离线积分或线性 Lyapunov 微分方程计算，在线控制只需读取当前误差和历史误差。

结合式 \eqref{eq:computed_torque} 和式 \eqref{eq:pdf_error_input}，理想 PDF 关节力矩控制律可写为
\begin{equation}
\boxed{
\tau
=
M_e(q_m)
\left[
\ddot q_d
+
K_0z(t)
-
K_c(t)z(t-h)
\right]
+
C_e(q_m,\dot q_m)\dot q_m
-
\hat d(t).
}
\label{eq:final_controller}
\end{equation}

在仿真实现中，为避免在示例程序中引入完整的离线矩阵增益求解，本文采用与二阶误差通道匹配的工程化平滑延迟项
\[
K_0=[-K_p\ -K_d],\qquad
-K_c(t)=\rho R_h(t)[-K_{ph}\ -K_{dh}],
\]
其中 \(K_p,K_d,K_{ph},K_{dh}\) 为正定对角矩阵，\(\rho>0\) 为延迟反馈幅值系数。该形式保留 PDF 的等待--作用周期结构，但并不等价于严格的 monodromy matrix 零化增益。因此，仿真结果用于验证控制结构和有界扰动下的跟踪性能；严格固定时间精确归零结论仍以式 \eqref{eq:kc_pdf} 的离线 PDF 增益为前提。

若考虑有界扰动，可引入饱和型鲁棒补偿
\[
\hat d(t)=k_d\,\operatorname{sat}(s/\varphi),
\qquad
s=\dot e+\Lambda e,
\]
其中 \(k_d>0\)、\(\varphi>0\)、\(\Lambda=\Lambda^\top>0\)。需要指出，加入非线性鲁棒补偿后，闭环不再是文献\cite{zhouFixedtimeStabilizationLinear2022}中的纯线性延迟系统，严格固定时间精确归零结论需要单独证明；工程上通常得到对扰动鲁棒的实际固定时间有界结果。
